% ============================================================
% Alpesh Sheth — Curriculum Vitae
% XeLaTeX · Latin Modern · dotted rules echoing asheth-git.github.io
% ============================================================
\documentclass[10pt,a4paper]{article}

\usepackage[a4paper,margin=1.65cm,bottom=1.8cm]{geometry}
\usepackage{fontspec}
\setmainfont{Latin Modern Roman}
\setmonofont{Latin Modern Mono}

\usepackage{xcolor}
\definecolor{ink}{HTML}{1B2430}
\definecolor{dim}{HTML}{55677A}
\definecolor{teal}{HTML}{0B7A71}
\definecolor{amber}{HTML}{C96F1A}
\definecolor{ruledot}{HTML}{9AB0C2}

\usepackage{tikz}
\usepackage{enumitem}
\usepackage{titlesec}
\usepackage[hidelinks]{hyperref}
\hypersetup{colorlinks=true, urlcolor=teal, linkcolor=teal}

\pagestyle{empty}
\color{ink}
\setlength{\parindent}{0pt}

% dotted rule, echoing the dot-matrix motif of the website
\newcommand{\dotrule}{%
  \tikz{\draw[color=ruledot, line width=1.5pt, line cap=round,
        dash pattern=on 0pt off 6pt] (0,0) -- (\linewidth,0);}}

\titleformat{\section}
  {\normalfont\ttfamily\small\bfseries\color{teal}\uppercase}
  {}{0em}{}[\vspace{2pt}\dotrule]
\titlespacing*{\section}{0pt}{14pt}{8pt}

% experience entry: role, institution, dates
\newcommand{\entry}[3]{%
  \textbf{#1}\hfill{\ttfamily\footnotesize\color{dim}#3}\\
  {\itshape\color{dim}#2}\vspace{2pt}}

\setlist[itemize]{leftmargin=1.1em, itemsep=1.5pt, topsep=3pt,
  label={\textcolor{amber}{\footnotesize\textbullet}}}

\begin{document}

% ------------------------------------------------------------ header
{\fontsize{27}{30}\selectfont Alpesh Sheth}\hfill
\begin{tabular}[b]{@{}r@{}}
  {\ttfamily\footnotesize \href{mailto:alpeshsheth.phy@gmail.com}{alpeshsheth.phy@gmail.com}}\\[1pt]
  {\ttfamily\footnotesize \href{https://www.linkedin.com/in/alpesh-sheth/}{linkedin.com/in/alpesh-sheth}}\\[1pt]
  {\ttfamily\footnotesize \href{https://asheth-git.github.io}{asheth-git.github.io}}
\end{tabular}

\vspace{4pt}
{\ttfamily\small\color{amber}PhD · Computational Condensed Matter Physics}

\vspace{6pt}
\dotrule

\vspace{8pt}
Computational physicist working on the electronic structure of quantum materials.
I construct effective multi-orbital models for correlated-electron systems, compute them
with many-body and first-principles methods on high-performance computing infrastructure,
and transfer the same modeling discipline to applied problems in machine learning and
optimization.

% ------------------------------------------------------------ experience
\section{Experience}

\entry{Doctoral Researcher}
      {Laboratoire Ondes et Matière d'Aquitaine (UMR 5798, CNRS), Université de Bordeaux}
      {Dec 2021–Nov 2025}
\begin{itemize}
  \item Constructed effective multi-orbital models separating the electronic physics of
        infinite-layer nickelates (RNiO\textsubscript{2}) from high-T\textsubscript{c} cuprates.
  \item Identified robust antiferromagnetic and charge-order tendencies driven by
        Fermi-surface nesting.
  \item Developed and optimized Python and Fortran modules for quantum-material simulation;
        workflow automation on high-performance computing clusters reduced compute time by
        approximately 35\,\%.
  \item Designed scalable frameworks to validate theoretical predictions against
        experimental data for complex material properties.
\end{itemize}

\entry{Research Intern, Data Science}
      {Uresearcher}
      {Jun 2022–Nov 2022}
\begin{itemize}
  \item Automated data-extraction pipelines through public interfaces, reducing manual
        pre-processing by approximately 70\,\%.
  \item Applied predictive machine-learning models to drug-design datasets, improving
        hit-prediction accuracy by approximately 15\,\%.
  \item Mined open-source molecular databases to identify candidate molecules for
        targeted applications.
\end{itemize}

\entry{Research Fellow}
      {The Maharaja Sayajirao University of Baroda}
      {Apr 2020–Nov 2021}
\begin{itemize}
  \item Performed density functional theory and molecular dynamics simulations of
        boron-based nanomaterials for radiation shielding.
  \item Designed modular high-performance computing workflows across more than twenty
        structural parameter sets, accelerating design cycles by approximately 50\,\%.
  \item Presented a density-functional study of twisted graphene and hexagonal boron
        nitride moiré heterostructures at ICAFMOD 2020.
\end{itemize}

\entry{Research Intern}
      {Physical Research Laboratory}
      {Jun 2018–Nov 2021}
\begin{itemize}
  \item Computed the non-linear alternating-current admittance of materials, improving
        predictive reliability by approximately 40\,\% over prior approaches.
  \item Analyzed transport properties through the density of states using second-order
        perturbation theory and the Kubo formalism.
  \item Determined critical size thresholds for superconducting nanoparticles used in
        magnetic sensor design.
\end{itemize}

% ------------------------------------------------------------ education
\section{Education}

\entry{PhD in Material Science}{Université de Bordeaux}{2021–2025}

\vspace{5pt}
\entry{MSc in Condensed Matter Physics}{The Maharaja Sayajirao University of Baroda}{2018–2020}

\vspace{5pt}
\entry{BSc in Physics}{The Maharaja Sayajirao University of Baroda}{2015–2018}

% ------------------------------------------------------------ publications
\section{Publications}

\begin{itemize}
  \item \textbf{A. Sheth} et al., \emph{Emergence and tunability of Fermi-pocket and
        electronic instabilities in layered Nickelates},
        Journal of Physics: Condensed Matter \textbf{37}, 125703 (2025).
        \href{https://doi.org/10.1088/1361-648X/ada908}{doi:10.1088/1361-648X/ada908}
  \item \textbf{A. Sheth}, \emph{Theoretical modelling of infinite-layer nickelates},
        PhD thesis, Université de Bordeaux (2025).
        \href{https://theses.hal.science/tel-05502836}{theses.hal.science/tel-05502836}
  \item B. Devanarayanan, A. Sharma, \textbf{A. Sheth} et al.,
        \emph{Background for the Self Consistent Renormalisation (SCR) Theory},
        \href{https://arxiv.org/abs/2108.12683}{arXiv:2108.12683} (2021).
\end{itemize}

% ------------------------------------------------------------ talks
\section{Invited Talks}

\begin{itemize}
  \item \emph{Modelling oxygen and rare-earth atomic impurities in nickelates},
        IIP-Natal WE-Heraeus Workshop on New Frontiers in Emergent Materials (2023).
        \href{https://www.youtube.com/live/tUv6uxbzKBw}{Recording available}.
\end{itemize}

% ------------------------------------------------------------ projects
\section{Selected Projects}

\begin{itemize}
  \item \textbf{Quantum Portfolio.} Full-stack application for portfolio optimization
        mapped onto the Ising model: live market data ingestion, constraint penalties for
        sector exposure and cardinality, a tunable risk-aversion parameter, and simulated
        annealing with variational free-energy tracking.
  \item \textbf{Graphene heterostructures.} Electronic properties of twisted graphene and
        hexagonal boron nitride bilayers from density functional theory, including moiré
        superlattice effects.
\end{itemize}

% ------------------------------------------------------------ skills
\section{Technical Skills}

\begin{tabular}{@{}p{4.2cm}p{12.2cm}@{}}
  {\ttfamily\footnotesize\color{dim}Languages} &
    Python, C++, Fortran, Julia, Bash, MATLAB, \LaTeX\\[3pt]
  {\ttfamily\footnotesize\color{dim}Many-body methods} &
    Green's functions, dynamical mean-field theory, exact diagonalization,
    second-order perturbation theory, Kubo formalism\\[3pt]
  {\ttfamily\footnotesize\color{dim}Simulation packages} &
    VASP, Quantum ESPRESSO, LAMMPS, CP2K, TRIQS\\[3pt]
  {\ttfamily\footnotesize\color{dim}Data and machine learning} &
    predictive modeling, statistical analysis, automated data pipelines\\[3pt]
  {\ttfamily\footnotesize\color{dim}Infrastructure} &
    high-performance computing (Slurm), Git, workflow automation
\end{tabular}

\end{document}
